\section{Theorem and assumptions}

We consider solutions to the linear Schrödinger equation
\[
i\partial_t u = -\Delta u + V(x)u,
\]
where \(u : \mathbb{R}^n \times \mathbb{R} \to \mathbb{C}\) and \(V\) is a suitable potential.

A representative unique continuation result (after Escauriaza–Kenig–Ponce–Vega) has the following form.

\begin{theorem}[Two-time Gaussian decay implies triviality]
Let \(u\) be a sufficiently regular solution to the Schrödinger equation above.
Assume that for two distinct times \(t_1 \neq t_2\),
\[
|u(x,t_1)| \le C e^{-a|x|^2}, \quad
|u(x,t_2)| \le C e^{-a|x|^2},
\]
for some constants \(C > 0\) and sufficiently large \(a > 0\).

Then, under appropriate conditions on \(V\),
\[
u(x,t) \equiv 0.
\]
\end{theorem}

\begin{remark}
The precise hypotheses depend on the dimension, regularity class, and assumptions on the potential \(V\).
We do not attempt to optimize constants or assumptions here, as the focus is structural.
\end{remark}
